\section{The sharp conductor bound for a fixed Gauss valuation}
\label{sec:fixed-radius}

Fix a prime $p$. Normalize the valuation of $\mathbf C_p$ by $v(p)=1$ and
write $v(0)=+\infty$. Set $e_h=(p-1)p^{h-1}$, $r_h=1/e_h$, and
$U_h=\{\zeta:\operatorname{ord}(\zeta)=p^h\}$. For a nonzero polynomial
$g(T)=\sum_n a_nT^n$ with coefficients in $\mathcal O_{\mathbf C_p}$ let
$G_\rho(g)=\min_n(v(a_n)+n\rho)$. The same definition applies to integral
formal series at every $\rho>0$. A shell $h$ is \emph{blocked} at tolerance
$\delta\ge0$ when every $\zeta\in U_h$ satisfies
$v(g(\zeta-1))>G_\rho(g)+\delta$. Equality is a successful specialization.

\paragraph{Theorem (exact worst conductor).}
Fix $h_0\ge1$ and $\rho\ge r_{h_0}$. The smallest value of $G_\rho(g)$
among all integral polynomials blocking the $L$ consecutive shells starting
at $h_0$ at tolerance zero is
\[
                        Lp^{1-h_0}.
\]
Consequently, given $G_\rho(g)\le B$, the least endpoint that guarantees
a successful specialization in $[h_0,H]$ for every $g$ is
\[
                 \boxed{H=h_0+\lfloor p^{h_0-1}B\rfloor.}
\]
Both assertions hold for integral formal series. The same endpoint works
for a series in any finite number of variables, with independent choices
of coordinates on the indicated shells.

The prior \emph{Exact torsion specialization} report gave the sufficient
endpoint $h_0+\lceil p^{h_0}B(h_0+B+2)\rceil$ and left its optimal value
open. The prior \emph{Torsion and zero data} report computed the extremal
loss \emph{within} a shell. Across shells the fixed Gauss valuation
retains the savings from roots on every earlier shell. Keeping these
savings yields an exact balance and the linear answer above.

\subsection{A calculation on one shell}
With a primitive $p^h$th root $\xi$, write $z_a=\xi^a-1$ for unit
residues $a\bmod p^h$. Let $D_h(m)$ be the largest possible value of
\[
   \min_a \sum_{i=1}^m \big(v(z_a-\alpha_i)-r_h\big),
        \qquad v(\alpha_i)=r_h.
\]
The following formula is the shell calculation in the earlier
\emph{Torsion and zero data} report:
\[
 D_h(m)=
 \begin{cases}
  e_h^{-1}\displaystyle\sum_{j=1}^{h-1}e_j\lfloor m/e_j\rfloor,
                    &0\le m<e_h,\\
  +\infty,&m\ge e_h.
 \end{cases}
\]
Here $v(z_a)=r_h$ follows by evaluating $\Phi_{p^h}(1)=p$, since
primitive conjugates yield factors of equal valuation. For each $\alpha_i$,
replace it by a shell point $z_b$ nearest to it. For every $a$,
$v(z_a-z_b)\ge\min(v(z_a-\alpha_i),v(\alpha_i-z_b))
=v(z_a-\alpha_i)$, so the replacement can only raise the minimum.

Now $m$ marks occupy unit residues. Write $N_j(a)$ for the number
congruent to $a$ modulo $p^j$. If $a$ is unmarked, then
$v(z_a-z_b)=p^{v_p(a-b)}/e_h$. Expanding $p^t-1=\sum_{j=1}^t e_j$
shows the excess at $a$ is $e_h^{-1}\sum_{j<h} e_jN_j(a)$.
Follow a least-occupied unit class modulo $p$ and then a least-occupied
child at each subsequent level. This gives
$N_j(a)\le\lfloor m/e_j\rfloor$ simultaneously, and ends at an unmarked
point if $m<e_h$. This proves the upper bound. Conversely, the first $m$
positive integers prime to $p$ give at least $\lfloor m/e_j\rfloor$
marks in every unit class modulo $p^j$, since each successive block of
$e_j$ unit exponents covers every such class. This attains the bound.
For $m\ge e_h$, covering the entire shell gives infinite excess.

\subsection{The exact budget for any finite set of shells}
In fact let $S$ be any nonempty finite set of positive indices with
$\rho\ge\max_{h\in S}r_h$. In increasing order of $h\in S$, define
\[
  m_h=\min\left\{0\le m\le e_h:
  D_h(m)>\delta+\sum_{s\in S,\,s<h}m_s(r_s-r_h)\right\}.
\]
The exact minimum budget for blocking every shell in $S$ is then
\[
                           \sum_{h\in S}m_hr_h.                 \tag{*}
\]
To prove necessity, factor $g=b\prod_i(T-\alpha_i)$ in $\mathbf C_p$.
Gauss multiplicativity gives
$G_u(g)=v(b)+\sum_i\min(u,v(\alpha_i))$. Let $M_h$ count roots of
valuation $r_h$, with multiplicity. Then
$G_\rho(g)\ge G_\rho(g)-G_0(g)\ge\sum_{h\in S}M_hr_h$ by integrality.
At shell $h$ all other roots contribute exactly their Gauss term at $r_h$.
The savings relative to the target include
\[
 G_\rho(g)-G_{r_h}(g)\ \ge\sum_{s\in S,\,s<h} M_s(r_s-r_h).
\]
Thus blocking forces $D_h(M_h)>\delta+\sum_{s<h}M_s(r_s-r_h)$, with
$s$ restricted to $S$. Induction in $S$ gives $M_h\ge m_h$. The
argument permits roots at zero or of negative valuation.

For equality, take on each shell the first $m_h$ unit exponents in the
construction above, and form the monic polynomial with exactly these
shifted roots. Its Gauss valuation is $\sum m_hr_h$. At each selected
shell the savings equal the displayed sum exactly, and the minimum
excess over $G_{r_h}$ is $D_h(m_h)$. The strict inequality in the
recurrence blocks every shell. The coefficients belong to
$\mathbf Q_p(\mu_{p^{\max S}})$.

\subsection{Evaluation for consecutive shells and zero tolerance}
For $h=h_0+k-1$, put $c=p^{1-h_0}$, so $e_k/e_h=c$. We claim
$m_h=e_k$. Assume inductively the previous values have this form.
Their savings at shell $h$ are
\[
   d_h=(k-1)c-\frac{p^{k-1}-1}{e_h}.
\]
Indeed every earlier shell costs $c$ and their total number of marks
is $\sum_{i<k}e_i=p^{k-1}-1$. Direct evaluation gives
\[
  D_h(e_k-1)=
       \frac{(k-1)e_k-(p^{k-1}-1)}{e_h}=d_h.
\]
Here only $j<k$ contribute and $\lfloor(e_k-1)/e_j\rfloor=e_k/e_j-1$
for those $j$. The strict blocking test fails at $e_k-1$. If $k=h$ then
$D_h(e_k)=+\infty$. If $k<h$, the summand with $j=k$ appears, giving
$D_h(e_k)=ke_k/e_h>d_h$. This proves the claim from the first step,
including $p=2$ and $h_0=k=h=1$. Equation~$(*)$ gives $Lc$.
The forced number of roots is also sharp: its sum is
$\sum_{k=1}^L e_k=p^L-1$, achieved by the same monic polynomial.
For polynomials of degree at most $D$ and budget at most $B$, the exact
first guaranteed endpoint is therefore
\[
 h_0+\min\bigl(\lfloor p^{h_0-1}B\rfloor,
                    \lfloor\log_p(D+1)\rfloor\bigr).
\]
A budget $B$ can block precisely up to $\lfloor B/c\rfloor$
initial consecutive shells in the worst case. If this number is zero,
the first shell itself is successful. When $p=3$, $h_0=2$ and $B=1$,
the extremizer blocks levels $2,3,4$ with $2,6,18$ roots respectively.
Level $5$ is guaranteed, improving the previous sufficient level $47$.

\paragraph{Passing to power series and several variables.}
For a fixed finite set of shells and a nonzero integral series, truncate
after degree $N$ with $(N+1)\min_{h\in S} r_h>G_\rho(g)+\delta$.
The tail has larger valuation than the comparison threshold at every
test point, and the truncation preserves $G_\rho$. This transfers the
polynomial result to series, including the sharpness examples.

In several variables, choose a coefficient $a_{\mathbf n}$ attaining the
weighted Gauss minimum $A=\min_{\mathbf n}(v(a_{\mathbf n})+\rho|\mathbf
n|)$. It exists since the coefficients are integral and $\rho>0$.
Within the first variable take the series $b(T_1)$ obtained by fixing all
other indices at $(n_2,\ldots,n_d)$. Then $G_\rho(b)+
\rho\sum_{j>1}n_j=A$. The one-variable theorem selects $z_1$ with
$v(b(z_1))\le G_\rho(b)$, so the remaining Gauss minimum is at most
$A$. Successive evaluation and induction give the same endpoint
independently in every coordinate. A polynomial depending on one
coordinate already forces the stated worst endpoint. Applying this to
a finite product gives an aggregate valuation bound by Gauss
multiplicativity.

\paragraph{Relation to earlier maximum principles.}
Kriz and Nordentoft, \emph{Horizontal $p$-adic $L$-functions},
arXiv:2310.20678v3, Propositions 2.7 and 2.11, establish discrete maximum
principles for Fourier transforms of measures on finite abelian groups.
In the cyclic case their second proposition gives the one-shell threshold
$p^{1-h}$ relative to the trivial-character value. Here the reference
is the coefficient Gauss valuation at positive radius, and the proof
computes the obstruction from every previous shell and realizes the
sharp cost simultaneously. The optimality claim concerns the explicit
extremal problem above.   
